\section{Answer 2: testing connectedness from \texorpdfstring{$A$}{A} alone}\label{sec:answer2}
% claims used here: C9 C30 C31 C32 C33 C34 C35 C36 C37 C38 C39 m3-1 m3-3 m3-4 m3-5 m3-8 m3-9 m3-10 m3-11 m3-12 m3-13 au-3

Question~2 asks how to know from $A$ whether the graph is connected
(\Cref{def:connected}). The idea is to turn the question about trips into a
question about counts. By \cref{thm:walkcount}, the entries of the powers of $A$
count walks. By \cref{lem:walkpath}, a walk between two vertices contains a path
between them, and by \cref{cor:bound} that path is short. So it is enough to
look at walks of length at most $N-1$.

\subsection{The running total \texorpdfstring{$S$}{S}}\label{sec:S}

\begin{plainwords}
We need a name for the running total. It adds the counts of trips with one road,
with two roads, and so on. We stop one step short of the number of houses. The
whole check uses this one total.
\end{plainwords}

\begin{definition}[The matrix $S$; new notation, not from the textbook]\label{def:S}
For a graph with $N\ge2$ vertices, let
\begin{equation}\label{eq:S}
S=A+A^2+\dots+A^{N-1},
\end{equation}
the sum of the first $N-1$ powers of $A$, added entry by entry.
\end{definition}

\begin{remark}\label{rem:S}
By \cref{thm:walkcount}, $S_{ij}$ is the number of walks from $n_i$ to $n_j$ whose length is between $1$ and $N-1$. Hence every entry of $S$ is a non-negative integer. The diagonal of $S$ counts closed walks and plays no part in the test (\cref{rem:diagonal}). \kindtag{derived here}
\end{remark}

\begin{example}\label{ex:S}
For the board graph $B$, $N=3$, so $S=A+A^2$:
\[
S=\MBoardA+\MBoardAsq=\MBoardS .
\]
\end{example}

\begin{nonexample}\label{nex:S}
For $N=3$, neither $A$ alone nor $A^2$ alone is $S$. The matrix $A^2$ of $B$ has the zero entry $(1,2)$ although $n_1$ and $n_2$ are joined by an edge. The matrix $A$ has the zero entry $(1,3)$ although $n_1$ and $n_3$ are connected through $n_2$. Either one alone would wrongly report ``not connected''. \kindtag{derived here}
\end{nonexample}

\subsection{The test and its proof}\label{sec:test}

\begin{plainwords}
We want a quick check for a village. We only have the table of shared roads. From
the table we can count trips. Count the trips with one road, with two roads, and
so on. Add the counts for each pair of houses. If no pair of different houses has
a total of zero, the village is connected. We never need trips longer than one
road fewer than the number of houses. A trip that wastes time can always be
shortened. So longer trips tell us nothing new.
\end{plainwords}

\begin{proposition}[One pair at a time]\label{prop:pair}
Let $G$ have $N\ge2$ vertices and let $i\ne j$. The following are equivalent.
\begin{enumerate}
\item[(i)] There is a path from $n_i$ to $n_j$ (distinct vertices; by \cref{prop:readings} this is the same as asking for a trail or a walk).
\item[(ii)] There is an integer $k$ with $1\le k\le N-1$ and $(A^k)_{ij}>0$.
\item[(iii)] $S_{ij}>0$.
\end{enumerate}
\end{proposition}

\begin{proof}
(i)$\Rightarrow$(ii). Let $P$ be a path from $n_i$ to $n_j$, of length $k$. By \cref{cor:bound}, $k\le N-1$. Because $n_i\ne n_j$, a walk of length $0$ cannot join them, so $k\ge1$. $P$ is a walk of length $k$ from $n_i$ to $n_j$, so by \cref{thm:walkcount} the number $(A^k)_{ij}$ is at least $1$.

(ii)$\Rightarrow$(iii). $S_{ij}=A_{ij}+(A^2)_{ij}+\dots+(A^{N-1})_{ij}$. Each term is a count, hence a non-negative integer (\cref{cor:walkcount}(a)). A sum of non-negative numbers is at least as large as any one of its terms. The term with the index $k$ of (ii) is positive, so $S_{ij}>0$.

(iii)$\Rightarrow$(ii). If $S_{ij}>0$, a sum of non-negative terms is positive, so at least one term $(A^k)_{ij}$ with $1\le k\le N-1$ is positive.

(ii)$\Rightarrow$(i). If $(A^k)_{ij}>0$ with $k\ge1$, then by \cref{thm:walkcount} there is at least one walk of length $k$ from $n_i$ to $n_j$. Since $n_i\ne n_j$, \cref{lem:walkpath} gives a path from $n_i$ to $n_j$.
\end{proof}

\begin{plainwords}
Here is the answer in short. Pick any two different houses. Look for a trip
between them. It never needs more roads than the number of houses, less one. The
village is connected when every such pair has one. The table of counts lets us
read this off.
\end{plainwords}

\begin{theorem}[The connectivity test]\label{thm:test}
Let $G$ be a graph with $N\ge2$ vertices and adjacency matrix $A$. Then $G$ is connected if and only if every entry off the diagonal of $S=A+A^2+\dots+A^{N-1}$ is positive.
\end{theorem}

The statement appears in Paramadevan and Sotheeswaran's Theorem~2~\cite{paramadevan2021adjacency} and in Weisstein's entry on connected graphs~\cite{weisstein_connectedgraph}. The proof in the first is incomplete (\cref{sec:provenance} quotes it), so the paper proves the test itself and cites these sources for the statement only. The textbook does not state it in the pages checked \kindtag{qualified}. The proof is \kindtag{derived here}.

\begin{remark}\label{rem:test}
The test allows any graph with $N\ge2$ vertices, connected or not. It looks only at entries with $i\ne j$. It forbids requiring positive diagonal entries (\cref{rem:diagonal}) and forbids stopping at $A^{N-2}$ (\cref{prop:sharp}).
\end{remark}

\begin{example}\label{ex:test}
The board graph $B$ has $S$ as in \cref{ex:S}: no zero entry, so $B$ is connected, which agrees with \cref{ex:connected}.
\end{example}

\begin{nonexample}\label{nex:test}
The graph $D$ (one edge $12$) has the matrix $S$ written out in \cref{sec:example}, with zeros at $(1,3)$, $(2,3)$, $(3,1)$, $(3,2)$. So $D$ is not connected, which agrees with \cref{nex:connected}.
\end{nonexample}

\begin{proof}
By \cref{def:connected} and \cref{prop:readings}, $G$ is connected if and only if every two vertices are joined by a path. For a pair with $u=v$ this is always true, by the one-vertex path (\cref{rem:connected}). For a pair $n_i,n_j$ with $i\ne j$ it is statement~(i) of \cref{prop:pair}. By \cref{prop:pair} that is equivalent to $S_{ij}>0$. So $G$ is connected if and only if $S_{ij}>0$ for every $i\ne j$.
\end{proof}

\paragraph{Why the argument is valid.}
It uses three results and nothing else: the walk count (\cref{thm:walkcount}), the walk-to-path lemma (\cref{lem:walkpath}) with its length bound (\cref{cor:bound}), and the fact that a count is a non-negative integer, so a sum is positive exactly when a term is. The last step would fail with subtraction or arithmetic modulo a number, which can cancel positive terms. The exponent $N-1$ is the largest one needed because a path has at most $N-1$ edges; adding higher powers cannot change which off-diagonal entries are positive.

\begin{plainwords}
The same counts also tell us the pieces of a village. Two different houses are in
the same piece exactly when some trip joins them. To find the distance, count the
trips of one road, then of two roads, and so on. The first number of roads that has
a trip is the distance.
\end{plainwords}

\begin{corollary}[Components and distances from the powers]\label{cor:components}
For any graph with $N\ge2$ vertices, connected or not, and $i\ne j$: (a)~$S_{ij}>0$ if and only if $n_i$ and $n_j$ lie in the same component. (b)~If they do, $d(n_i,n_j)$ is the smallest $k\ge1$ with $(A^k)_{ij}>0$; if they do not, no power has a positive $(i,j)$ entry and $d(n_i,n_j)=+\infty$. \kindtag{derived here}
\end{corollary}

\begin{proof}
(a) \Cref{prop:pair} never uses that $G$ is connected. Its statement~(i) says $n_i\sim n_j$, which is the definition of lying in one class, that is in one component (\cref{def:component}).
(b) Let $k_0$ be the smallest $k\ge1$ with $(A^k)_{ij}>0$. By \cref{thm:walkcount} a walk of length $k_0$ exists and no shorter walk does. By \cref{lem:walkpath} a path of length at most $k_0$ exists; being a walk, it has length at least $k_0$. So the shortest walk has length $k_0$, and the distance equals the length of the shortest walk (\cref{app:distance}). If no power has a positive entry, there is no walk and the distance is $+\infty$.
\end{proof}

The distance statement is also made by Lehman, Leighton and Meyer, and by Paramadevan and Sotheeswaran in their Theorem~3~\cite{lehman_mcs_adjacency,paramadevan2021adjacency}. The entry $S_{ij}$ tells whether $n_i$ and $n_j$ are connected, not which path joins them.

\begin{remark}[The diagonal, and the cases $N=1$ and $N=2$]\label{rem:diagonal}
For $N=2$ and one edge, $S=A$ has zeros on the diagonal although the graph is connected (by script: $A=\MTwoVertexA$). So the diagonal of $S$ cannot be required to be positive. For $N\ge3$ and $G$ connected the diagonal of $S$ is positive anyway, because every vertex has a neighbour (an isolated vertex would have no path to the other $N-1$ vertices) and $(A^2)_{ii}=\deg(n_i)$ (\cref{cor:walkcount}(d)); the test does not need this. For $N=1$ there are no pairs of different vertices and the graph is connected; the sum $S$ is empty, so the test is stated for $N\ge2$. \kindtag{derived here}
\end{remark}

\subsection{The test as a procedure}\label{sec:procedure}

\begin{plainwords}
Here is the check, step by step. First write the table of the village. Then
multiply the table by itself again and again. After each step, add the new counts
to a running total. Stop after one step fewer than the number of houses. Look at
every pair of different houses. If any total is zero, those two houses cannot
reach each other. If no total is zero, everyone can reach everyone.
\end{plainwords}

\begin{test}[Matrix test for connectedness]\label{test:matrix}
\emph{Input:} the adjacency matrix $A$ of a graph with $N\ge2$ vertices.
\begin{enumerate}
\item Set $P=A$ and $S=A$.
\item For $k=2,3,\dots,N-1$: replace $P$ by $PA$ (now $P=A^k$) and add $P$ to $S$.
\item Look at the entries $S_{ij}$ with $i\ne j$. If one of them is $0$, report ``not connected''. If none is $0$, report ``connected''.
\end{enumerate}
\end{test}

\begin{remark}\label{rem:procedure}
Correctness is \cref{thm:test}. The loop makes $N-2$ matrix products, the last one producing $A^{N-1}$. A zero at position $(i,j)$ means by \cref{cor:components} that $n_i$ and $n_j$ are in different components, so the procedure also lists the pieces. The test forbids stopping early: it needs every power up to $N-1$ (\cref{prop:sharp}).
\end{remark}

\begin{example}\label{ex:procedure}
On $B$ ($N=3$) the loop runs once ($k=2$): $P=A^2$, $S=A+A^2$, no off-diagonal zero, so ``connected''. The matrices are in \cref{sec:example}.
\end{example}

\begin{nonexample}\label{nex:procedure}
Reporting after $k=1$ only (the table $A$ itself) is not the test: on $B$ it would report ``not connected'' because $A_{13}=0$.
\end{nonexample}

\subsection{An equivalent form with one table}\label{sec:iplusa}

\begin{plainwords}
There is a second form of the check that uses one table only. Let a trip wait at
a house for a step. Then count trips of one fixed length. Waiting lets short
trips count too. The village is connected exactly when no pair of houses has a
count of zero.
\end{plainwords}

The table that makes a trip wait is the identity matrix $I$, which has ones on the diagonal and zeros elsewhere: ``stay at the same house'' is one allowed step, just as ``drive one road'' is the step given by $A$. When $(I+A)$ is multiplied by itself $m$ times, each of the $m$ factors offers a choice between waiting ($I$) and driving one road ($A$). A choice with $k$ drives and $m-k$ waits gives the product $A^k$ (because $IX=X$ for every matrix $X$), and there are $\binom{m}{k}$ such choices. So
\begin{equation}\label{eq:binomial}
(I+A)^m=\binom{m}{0}I+\binom{m}{1}A+\binom{m}{2}A^2+\dots+\binom{m}{m}A^m ,
\end{equation}
with $A^0=I$. Every coefficient is a positive integer. \Cref{app:iplusa} gives the argument in detail.

\begin{plainwords}
So one table is enough. Take the table that allows waiting and raise it to a high
enough power. A zero anywhere means that two houses cannot reach each other.
Otherwise every house can reach every other house.
\end{plainwords}

\begin{proposition}[The form with $I+A$]\label{prop:iplusa}
Let $G$ have $N\ge2$ vertices, let $I$ be the $N\times N$ identity matrix, and let $m\ge N-1$ be an integer. Then (a)~for $i\ne j$, $\bigl((I+A)^m\bigr)_{ij}>0$ if and only if $S_{ij}>0$; (b)~every diagonal entry of $(I+A)^m$ is at least $1$. Hence $G$ is connected if and only if $(I+A)^{N-1}$ has no zero entry. \kindtag{derived here; proof in \cref{app:iplusa}}
\end{proposition}

Bai states the case $m=N-1$ as his Theorem~2, with a proof sketch that counts ``paths'' where walks are meant~\cite{bai2025adjacency}. That statement is therefore \kindtag{qualified}; the proof in the appendix replaces the sketch. For the board graph and for $D$ the matrices $(I+A)^2$ are in \cref{sec:example}, and the zeros fall in the same places as the zeros of $S$.

\subsection{Why \texorpdfstring{$N-1$}{N-1} and not less}\label{sec:sharp}

\begin{plainwords}
Take four houses in a row. To go from the first house to the last you must pass
every house in between. That takes three roads. If we counted only trips of up to
two roads, we would miss this pair. We would say the village is cut in two, and
we would be wrong. The same idea works for any row of houses.
\end{plainwords}

\begin{proposition}[$N-1$ cannot be lowered to $N-2$]\label{prop:sharp}
For every $N\ge3$ there is a connected graph with $N$ vertices for which the sum $A+\dots+A^{N-2}$ has a zero entry off the diagonal. \kindtag{derived here}
\end{proposition}

\begin{proof}
Take the graph with edges $1\text{--}2,\ 2\text{--}3,\ \dots,\ (N-1)\text{--}N$, a row of $N$ houses; it is connected. Each step of a walk changes the vertex number by exactly $1$. To go from vertex $1$ to vertex $N$ the number must increase by $N-1$, so at least $N-1$ steps are needed. A walk of length $k<N-1$ from $1$ to $N$ does not exist, so by \cref{thm:walkcount} $(A^k)_{1N}=0$ for $1\le k\le N-2$. Hence the $(1,N)$ entry of $A+\dots+A^{N-2}$ is $0$, although the graph is connected.
\end{proof}

For $N=3$ the graph is the board graph and the sum $A+\dots+A^{N-2}$ is $A$ itself, with $A_{13}=0$. For $N=4$ (\cref{fig:path4}) every matrix is written out below; the sum of the first $N-2=2$ powers has zeros at $(1,4)$ and $(4,1)$, and the sum of the first $N-1=3$ powers has none.

\begin{figure}[ht]
\centering
\setlength{\unitlength}{1pt}
\begin{picture}(190,40)(0,-5)
  \put(10,15){\circle*{6}} \put(65,15){\circle*{6}} \put(120,15){\circle*{6}} \put(175,15){\circle*{6}}
  \put(10,15){\line(1,0){55}} \put(65,15){\line(1,0){55}} \put(120,15){\line(1,0){55}}
  \put(7,22){$1$} \put(62,22){$2$} \put(117,22){$3$} \put(172,22){$4$}
\end{picture}

\caption{Four houses in a row, the graph used in \cref{ex:sharp}.}
\label{fig:path4}
\end{figure}

\begin{example}[The bound $N-1$ is needed; four houses in a row]\label{ex:sharp}
With $N=4$, the matrices from the script are
{\small\begin{align*}
A&=\MPathA, & A^2&=\MPathAsq,\\
A^3&=\MPathAcube, & A+A^2&=\MPathSshort,\\
A+A^2+A^3&=\MPathSfull .
\end{align*}}
With $N-2=2$ powers the $(1,4)$ entry is $0$, so the test would wrongly say ``not connected''. With $N-1=3$ powers every entry is positive.
\end{example}

\begin{table}[ht]
\centering
\caption{The row of $N$ houses, $N=3,\dots,8$ (by script). The first power with a non-zero $(1,N)$ entry is always $N-1$.}
\label{tab:sharp}
\small
\begin{tabular}{@{}cp{3.4cm}p{3.8cm}p{3.8cm}@{}}
\toprule
$N$ & smallest $k$ with $(A^k)_{1N}>0$ & $(1,N)$ entry of $A+\dots+A^{N-2}$ & $(1,N)$ entry of $A+\dots+A^{N-1}$ \\
\midrule
\TabPathSharpBody
\bottomrule
\end{tabular}
\end{table}

Using only the single power $A^{N-1}$ is also wrong: on the board graph $A^2$ has the zero entry $(1,2)$.

\begin{remark}[Which graphs the $N-2$ version gets wrong]\label{rem:failcount}
The test with only $N-2$ powers is wrong on exactly the labelled rows of $N$ houses, that is, the path graphs with the vertices numbered in any order. There are $N!/2$ of them for $N\ge2$ ($N!$ orders along the row, each row counted twice because it can be read from either end): $1,3,12,60,360$ for $N=2,\dots,6$. These are the numbers in the last column of \cref{tab:exhaustive}, which the script counted by brute force. \kindtag{derived here} Reason: the $N-2$ version can only err by saying ``not connected'' for a connected graph, because a positive entry always comes from a real walk. By \cref{cor:components}(b) it says so exactly when some pair of vertices has distance $N-1$. A shortest path between two such vertices passes through all $N$ vertices, and a shortest path has no edge between vertices that are not neighbours on it, since that edge would give a shorter trip. So the graph is that row of houses. Conversely, \cref{prop:sharp} shows that a row fails. For $N=2$ the sum of $N-2=0$ powers is empty, and an empty sum is the zero matrix, so the test says ``not connected'' on the one connected graph.
\end{remark}

\subsection{Breadth-first search, the alternative}\label{sec:bfs}

\begin{plainwords}
There is another way, and it needs no table of counts. Stand at one house and
mark it. Then mark every house next to it. Then mark every house next to those.
Keep going until no new house appears. If every house is marked, the village is
connected.
\end{plainwords}

\emph{Breadth-first search} (BFS) from a start vertex $s$ is this procedure.
\begin{enumerate}
\item Mark $s$ and put it on a waiting list, called a \emph{queue}: vertices leave it in the order in which they entered.
\item Take the first vertex off the list. Mark every neighbour of it that is not yet marked (these vertices are now \emph{discovered}) and put them at the end of the list.
\item Repeat step~2 until the list is empty.
\item Report ``connected'' if every vertex is marked, and ``not connected'' otherwise.
\end{enumerate}
The correctness statement in the source is that for every vertex $v$ there is a walk from $s$ to $v$ if and only if $v$ is discovered when the search ends, so $G$ is connected if and only if every vertex is discovered when BFS starts from any single vertex $s$~\cite{schost2026bfs}. When the graph is stored as adjacency lists (for each vertex, the list of its neighbours), the source gives a running time at most a fixed multiple of $N+|E|$~\cite{schost2026bfs}. BFS and the matrix test agree on all labelled graphs with $2,3,4,5,6$ vertices ($2$, $8$, $64$, $1024$, $32768$ graphs; a check, not a proof, \cref{app:crosschecks}), and BFS also lists the components directly. The two methods are compared for cost in \cref{sec:discussion}.
